% Part: model-theory
% Chapter: models-of-arithmetic
% Section: introduction

\documentclass[../../../include/open-logic-section]{subfiles}

\begin{document}

\olfileid[hi]{mod}{mar}{int}
\section{परिचय}

अंकगणित का \emph{मानक प्रतिरूप} वह
!!{structure}~$\Struct{N}$ है जिसमें $\Domain{N} = \Nat$ है और
$\Obj{0}$, $\prime$, $+$, $\times$ तथा $<$ की व्याख्या वैसी ही की जाती
है जैसी अपेक्षित है। अर्थात् $\Obj{0}$ का मान $0$ है, $\prime$ उत्तरवर्ती
फलन है, $+$ की व्याख्या योग के रूप में और $\times$ की व्याख्या~$\Nat$ की
संख्याओं के गुणन के रूप में होती है। विशेषतः,
\begin{align*}
  \Assign{\Obj{0}}{N} & = 0\\
  \Assign{\prime}{N}(n) & = n + 1\\
  \Assign{+}{N}(n, m) & = n + m\\
  \Assign{\times}{N}(n, m) & = nm
\end{align*}
निस्संदेह, $\Lang{L_A}$ की ऐसी संरचनाएँ भी हैं जिनके प्रांत~$\Nat$ से
भिन्न हैं। उदाहरण के लिए, हम $\Struct{M}$ ले सकते हैं जिसका प्रांत
$\Domain{M} = \{a\}^*$ है (एकमात्र प्रतीक~$a$ के परिमित अनुक्रम, अर्थात्
$\emptyset$, $a$, $aa$, $aaa$, \dots), और जिसकी व्याख्याएँ हैं
\begin{align*}
  \Assign{\Obj{0}}{M} & = \emptyset\\
  \Assign{\prime}{M}(s) & = s \concat a\\
  \Assign{+}{M}(n, m) & = a^{n + m}\\
  \Assign{\times}{M}(n, m) & = a^{nm}
\end{align*}
ये दोनों संरचनाएँ इस अर्थ में “मूलतः एक जैसी” हैं कि अंतर केवल उनके
!!{domain}s के !!{element}s में है, इस बात में नहीं कि व्याख्या-फलनों द्वारा
उन !!{domain}s के !!{element}s को परस्पर कैसे संबंधित किया गया है। हम कहते
हैं कि ये दोनों !!{structure}s \emph{समरूपी} हैं।

संहतता प्रमेय का एक सरल परिणाम यह है कि~$\Struct{N}$ में सत्य किसी भी
सिद्धांत के ऐसे प्रतिरूप भी होते हैं जो~$\Struct{N}$ के समरूपी नहीं हैं।
ऐसी संरचनाएँ \emph{अमानक} कहलाती हैं। उनकी रोचक विशेषता यह है कि किसी मानक
प्रतिरूप के !!{element}s (अर्थात् $\Struct{N}$ के, और उसके समरूपी सभी
!!{structure}s के भी) मानक संख्यांकों~$\num{n}$ के मानों से पूरी तरह मिल
जाते हैं, अर्थात्
\[
\Domain{N} = \Setabs{\Value{\num{n}}{N}}{n \in \Nat}
\]
पर अमानक प्रतिरूपों में ऐसा नहीं होता: यदि $\Struct{M}$ अमानक है, तो कम-से-कम
एक $x \in \Domain{M}$ ऐसा है कि $x \neq \Value{\num{n}}{M}$ प्रत्येक~$n$ के
लिए।

ये अमानक अवयव काफ़ी दिलचस्प हैं: वे “अनंत प्राकृत संख्याएँ” हैं। किंतु उनका
अस्तित्व एक अर्थ में अपूर्णता की परिघटनाओं को भी स्पष्ट करता है। उदाहरण के
लिए पियानो अंकगणित का संगतता-कथन $\OCon[\Th{PA}]$ लें, अर्थात् $\lnot
\lexists[x][\OPrf[\Th{PA}](x, \gn{\lfalse})]$। चूँकि $\Th{PA}$ न तो
$\OCon[\Th{PA}]$ को सिद्ध करता है, न $\lnot \OCon[\Th{PA}]$ को, इसलिए
दोनों में से किसी को भी $\Th{PA}$ में संगत रूप से जोड़ा जा सकता है। चूँकि
$\Th{PA}$ संगत है, $\Sat{N}{\OCon[\Th{PA}]}$, और फलतः
$\Sat/{N}{\lnot
  \OCon[\Th{PA}]}$। अतः $\Struct{N}$, $\Th{PA}
\cup \{\lnot \OCon[\Th{PA}]\}$ का प्रतिरूप \emph{नहीं} है, और इसके
सभी प्रतिरूप अमानक होने चाहिए।
$\Th{PA} \cup \{\lnot \OCon[\Th{PA}]\}$ के प्रतिरूपों में कोई ऐसा
!!{element} अवश्य होना चाहिए
जो $\lexists[x][\OPrf[\Th{PA}](\gn{\lfalse})]$ को सत्य बनाने वाला साक्षी
बने, अर्थात्~$\Th{PA}$ से किसी विरोधाभास की !!a{derivation} की एक ग्योडेल
(G\"odel) संख्या। ऐसा !!a{element} मानक नहीं हो सकता---क्योंकि
$\Th{PA} \Proves \lnot
\OPrf[\Th{PA}](\num{n}, \gn{\lfalse})$ प्रत्येक~$n$ के लिए।

\end{document}
